\documentclass[11pt,letterpaper]{article}
% KJO_preamble.tex
% Shared LaTeX preamble for all KJO documentation documents.
% \input this file immediately after \documentclass{}.

% --- Encoding and fonts -------------------------------------------------------
\usepackage[utf8]{inputenc}
\usepackage[T1]{fontenc}
\usepackage{lmodern}
\usepackage{microtype}

% --- Page geometry ------------------------------------------------------------
\usepackage[
  letterpaper,
  top=1in, bottom=1in,
  left=1.25in, right=1in
]{geometry}

% --- Core utilities -----------------------------------------------------------
\usepackage{amsmath}
\usepackage{graphicx}
\usepackage{xcolor}
\usepackage{parskip}        % space between paragraphs; no indent
\usepackage{enumitem}
\usepackage{caption}
\usepackage{subcaption}
\usepackage{float}

% --- Tables -------------------------------------------------------------------
\usepackage{booktabs}
\usepackage{longtable}
\usepackage{array}
\usepackage{multirow}

% --- Code listings ------------------------------------------------------------
\usepackage{listings}
\lstset{
  basicstyle=\ttfamily\small,
  keywordstyle=\bfseries\color{KJOblue},
  commentstyle=\itshape\color{gray},
  stringstyle=\color{KJOgreen!70!black},
  breaklines=true,
  frame=single,
  framesep=4pt,
  xleftmargin=6pt,
  xrightmargin=4pt,
  numbers=left,
  numberstyle=\tiny\color{gray},
  numbersep=8pt,
  language=C++,
}

% --- TikZ / PGF ---------------------------------------------------------------
\usepackage{tikz}
\usetikzlibrary{
  arrows.meta,
  shapes.geometric,
  shapes.misc,
  shapes.symbols,
  positioning,
  fit,
  calc,
  backgrounds,
  decorations.pathreplacing,
  decorations.markings,
  matrix,
}
\usepackage{pgfplots}
\pgfplotsset{compat=1.18}

% --- Headers / footers --------------------------------------------------------
\usepackage{fancyhdr}
\usepackage{lastpage}
\pagestyle{fancy}
\fancyhf{}
\renewcommand{\headrulewidth}{0.4pt}
\renewcommand{\footrulewidth}{0.4pt}
% Per-document: \lhead{}, \rhead{}, \cfoot{} are set in the main .tex file.

% --- Section formatting -------------------------------------------------------
\usepackage{titlesec}
\titleformat{\section}{\large\bfseries\color{KJOblue}}{{\thesection}}{0.8em}{}[\titlerule]
\titleformat{\subsection}{\normalsize\bfseries\color{KJOblue!80!black}}{{\thesubsection}}{0.8em}{}
\titleformat{\subsubsection}{\normalsize\bfseries}{{\thesubsubsection}}{0.8em}{}

% --- Hyperlinks ---------------------------------------------------------------
\usepackage[
  colorlinks=true,
  linkcolor=KJOblue,
  urlcolor=KJOblue!70!black,
  citecolor=KJOblue,
  pdfborder={0 0 0},
]{hyperref}

% --- Color palette -----------------------------------------------------------
\definecolor{KJOblue}   {RGB}{ 30,  90, 160}   % RCU
\definecolor{KJOred}    {RGB}{180,  30,  30}   % EMU
\definecolor{KJOgreen}  {RGB}{ 20, 120,  60}   % GSEMU
\definecolor{KJOorange} {RGB}{200, 110,  20}   % GSE / pressurization
\definecolor{KJOgray}   {RGB}{100, 100, 100}   % Buses / links
\definecolor{KJOpurple} {RGB}{110,  40, 140}   % Shared libraries
\definecolor{KJObg}     {RGB}{248, 248, 252}   % Light background for code

% --- Convenience macros -------------------------------------------------------

% Hardware component name (small caps)
\newcommand{\hw}[1]{\textsc{#1}}

% Software module name (monospace)
\newcommand{\sw}[1]{\texttt{#1}}

% Command name (monospace bold)
\newcommand{\cmd}[1]{\texttt{\textbf{#1}}}

% Pin / signal name
\newcommand{\sig}[1]{\texttt{#1}}

% Note / callout box
\usepackage{tcolorbox}
\tcbuselibrary{skins}
\newtcolorbox{notebox}{
  enhanced,
  colback=KJOblue!5,
  colframe=KJOblue!40,
  fonttitle=\bfseries,
  title=Note,
  left=4pt, right=4pt, top=2pt, bottom=2pt,
}
\newtcolorbox{warningbox}{
  enhanced,
  colback=KJOred!5,
  colframe=KJOred!40,
  fonttitle=\bfseries\color{KJOred},
  title=Warning,
  left=4pt, right=4pt, top=2pt, bottom=2pt,
}

% Title page macro: \KJOtitlepage{title}{subtitle}{version}{date}{description}
\newcommand{\KJOtitlepage}[5]{%
  \begin{titlepage}
    \centering
    \vspace*{2cm}
    {\Huge\bfseries\color{KJOblue} #1 \par}
    \vspace{0.5cm}
    {\Large\color{KJOgray} #2 \par}
    \vspace{1.5cm}
    \rule{0.5\textwidth}{0.4pt}\par
    \vspace{1cm}
    {\large Ken Overton \par}
    \vspace{0.3cm}
    {\normalsize \textit{KJO Liquid Rocket Motor Control System} \par}
    \vspace{1.5cm}
    \begin{tabular}{ll}
      \textbf{Version:} & #3 \\[4pt]
      \textbf{Date:}    & #4 \\
    \end{tabular}
    \vspace{1.5cm}\par
    \begin{minipage}{0.75\textwidth}
      \centering
      \normalsize #5
    \end{minipage}
    \vfill
    {\small\color{KJOgray} \textcopyright\ 2025--2026 Ken Overton.
     All rights reserved. \par}
  \end{titlepage}%
}

\newcommand{\docversion}{0.1}
\newcommand{\docdate}{March 2026}
\lhead{\textsc{KJO Liquid Rocket Motor Control System}}
\rhead{\textsc{System Overview} --- v\docversion}
\cfoot{\thepage\ of \pageref{LastPage}}
\hypersetup{pdftitle={KJO System Overview}, pdfauthor={Ken Overton}}

\begin{document}
% KJO_tikz_styles.tex
% Shared TikZ node and edge styles for all KJO block diagrams.
% \input this file inside the document body before any tikzpicture environments.

\tikzset{
  % ── Hardware block (peripheral chip / PCB component) ──────────────────────
  hwblock/.style={
    rectangle, draw=KJOgray, line width=1pt,
    fill=KJOgray!10, rounded corners=4pt,
    minimum width=2.6cm, minimum height=0.8cm,
    font=\small, align=center, inner sep=4pt,
  },
  % Sub-variant: MCU (larger, blue border)
  mcublock/.style={
    hwblock,
    draw=KJOblue, fill=KJOblue!8,
    minimum width=3.0cm, minimum height=1.0cm,
    font=\small\bfseries,
  },
  % Sub-variant: power / battery
  powerblock/.style={
    hwblock,
    draw=KJOorange, fill=KJOorange!8,
  },
  % ── Software module box ────────────────────────────────────────────────────
  swmodule/.style={
    rectangle, draw=KJOblue, line width=0.8pt,
    fill=KJOblue!6, rounded corners=3pt,
    minimum width=2.8cm, minimum height=0.7cm,
    font=\small, align=center, inner sep=3pt,
  },
  % Sub-variant: shared library
  sharedlib/.style={
    swmodule,
    draw=KJOpurple, fill=KJOpurple!6,
    dashed,
  },
  % ── Bus / link annotation ──────────────────────────────────────────────────
  busline/.style={
    draw=KJOgray, line width=1.5pt,
    -{Stealth[length=6pt,width=4pt]},
  },
  bidbusline/.style={
    draw=KJOgray, line width=1.5pt,
    {Stealth[length=6pt,width=4pt]}-{Stealth[length=6pt,width=4pt]},
  },
  radioline/.style={
    draw=KJOblue!60, line width=1.5pt, dashed,
    {Stealth[length=6pt,width=4pt]}-{Stealth[length=6pt,width=4pt]},
  },
  gpioline/.style={
    draw=KJOgreen!60!black, line width=1.2pt,
    -{Stealth[length=5pt,width=4pt]},
  },
  % ── Flowchart shapes ──────────────────────────────────────────────────────
  flowproc/.style={
    rectangle, draw=KJOblue!70, line width=0.8pt,
    fill=KJOblue!8, rounded corners=3pt,
    minimum width=3.2cm, minimum height=0.7cm,
    font=\small, align=center, inner sep=3pt,
  },
  flowdec/.style={
    diamond, draw=KJOorange, line width=0.8pt,
    fill=KJOorange!8,
    minimum width=2.2cm, minimum height=0.9cm,
    font=\small, align=center, inner sep=1pt,
  },
  flowterminal/.style={
    rounded rectangle, draw=KJOblue, line width=1pt,
    fill=KJOblue!12,
    minimum width=2.8cm, minimum height=0.7cm,
    font=\small\bfseries, align=center, inner sep=3pt,
  },
  flowarrow/.style={
    draw=KJOgray, line width=0.9pt,
    -{Stealth[length=5pt,width=4pt]},
  },
  % ── Propellant schematic styles ────────────────────────────────────────────
  % Generic pipe line
  pipe/.style={
    draw=black, line width=1.8pt,
  },
  pipearrow/.style={
    pipe,
    -{Stealth[length=7pt,width=5pt]},
  },
  % Oxidizer (N2O) — blue
  oxpipe/.style={
    draw=KJOblue, line width=1.8pt,
  },
  oxpipearrow/.style={
    oxpipe,
    -{Stealth[length=7pt,width=5pt]},
  },
  % Fuel — red
  fuelpipe/.style={
    draw=KJOred, line width=1.8pt,
  },
  fuelpipearrow/.style={
    fuelpipe,
    -{Stealth[length=7pt,width=5pt]},
  },
  % GSE fill line — green
  gsepipe/.style={
    draw=KJOgreen!70!black, line width=1.8pt,
  },
  gsepipearrow/.style={
    gsepipe,
    -{Stealth[length=7pt,width=5pt]},
  },
  % Pressurization — orange dashed
  presspipe/.style={
    draw=KJOorange, line width=1.4pt, dashed,
  },
  presspipearrow/.style={
    presspipe,
    -{Stealth[length=6pt,width=4pt]},
  },
  % Tank (tall rectangle, rounded corners)
  tank/.style={
    rectangle, draw=black, line width=1.5pt,
    rounded corners=4pt, fill=white,
    minimum width=1.8cm, minimum height=2.4cm,
    font=\small\bfseries, align=center, inner sep=4pt,
  },
  n2otank/.style={
    tank, draw=KJOblue, fill=KJOblue!8,
  },
  fueltank/.style={
    tank, draw=KJOred, fill=KJOred!8,
  },
  srctank/.style={
    tank, draw=KJOgreen!70!black, fill=KJOgreen!6,
  },
  % Valve box (servo-actuated)
  valvebox/.style={
    rectangle, draw=black, line width=1pt,
    fill=white, minimum width=2.2cm, minimum height=0.65cm,
    font=\footnotesize\bfseries, align=center, inner sep=3pt,
  },
  gseval/.style={
    valvebox, draw=KJOgreen!70!black, fill=KJOgreen!8,
  },
  emuval/.style={
    valvebox, draw=KJOred, fill=KJOred!6,
  },
  % Component (generic inline component)
  component/.style={
    rectangle, draw=KJOgray, line width=0.8pt,
    fill=white!96!gray, minimum width=2.2cm, minimum height=0.6cm,
    font=\footnotesize, align=center, inner sep=3pt,
  },
  % Instrument callout circle (PT, TT, TC)
  instrument/.style={
    circle, draw=KJOgray, line width=0.8pt,
    fill=white, minimum size=0.6cm,
    font=\tiny\bfseries, align=center, inner sep=0pt,
  },
  % Sensor lead line
  sensorlead/.style={
    draw=KJOgray, line width=0.7pt,
  },
  % Boundary box (dashed outline for GSE / Rocket regions)
  gseboundary/.style={
    draw=KJOgreen!50, line width=1pt, dashed,
    rounded corners=6pt,
  },
  rocketboundary/.style={
    draw=KJOred!40, line width=1pt,
    rounded corners=6pt,
  },
  % Sequence diagram lifeline
  lifeline/.style={
    draw=KJOgray, line width=0.6pt, dashed,
  },
  % Sequence diagram message arrow
  seqarrow/.style={
    draw=KJOblue, line width=0.9pt,
    -{Stealth[length=5pt,width=4pt]},
  },
  % State node (FSM)
  statenode/.style={
    circle, draw=KJOblue, line width=1pt,
    fill=KJOblue!10,
    minimum size=1.2cm,
    font=\small\bfseries, align=center,
  },
  statetransition/.style={
    draw=KJOgray, line width=0.9pt,
    -{Stealth[length=5pt,width=4pt]},
    font=\footnotesize,
  },
}


\KJOtitlepage%
  {KJO Liquid Rocket Motor\\Control System}%
  {System Architecture and Operational Overview}%
  {\docversion}{\docdate}%
  {This document provides a high-level structural and functional overview of the
   KJO control system for a bi-propellant liquid rocket motor.  It describes the
   four electronic subsystems (RCU, EMU, GSEMU, LCMU), the propellant flow
   architecture, the inter-unit communication protocols, and the complete
   command set.  Detailed subsystem reference documents are listed in
   Section~\ref{sec:references}.}

\tableofcontents
\newpage

%═══════════════════════════════════════════════════════════════════════════════
\section{Introduction}
%═══════════════════════════════════════════════════════════════════════════════

The KJO system controls a bi-propellant liquid rocket motor using nitrous oxide
(N\textsubscript{2}O) as the oxidizer and a liquid fuel with regenerative
cooling.  Four autonomous microcontroller units share control responsibilities:

\begin{itemize}[nosep]
  \item \textbf{Remote Control Unit (RCU)} --- handheld operator terminal.
        Sends commands over a 900~MHz LoRa radio link; displays system status
        on a 3.5\textquotedbl\ colour touchscreen.
  \item \textbf{Engine Management Unit (EMU)} --- mounted in the rocket.
        Receives commands from the RCU, actuates the Oxidizer Feed and Fuel Feed
        servos, reads all sensors, and coordinates the Fill valve and load-cell
        readings with the GSEMU and LCMU over a shared CAN bus.
  \item \textbf{GSE Management Unit (GSEMU)} --- mounted on the launch pad.
        Controls the Fill valve servo and the umbilical quick-release servo;
        coordinates the fill-to-target sequence with the LCMU and reports
        status to the EMU, both over the shared CAN bus. A separate,
        ground-referenced disconnect-sense line (not CAN) gives both EMU and
        GSEMU instant, protocol-independent detection of umbilical separation.
  \item \textbf{Load Cell Management Unit (LCMU)} --- mounted on the test
        stand. Measures Ox load during oxidizer fill (supporting GSEMU's
        fill-to-target sequence) and thrust during test-stand burns, logging
        both to microSD; answers CAN queries from EMU and GSEMU.
\end{itemize}

RCU, EMU, and GSEMU are based on the Adafruit Feather~M0 (ATSAMD21G18A,
48~MHz ARM~Cortex-M0+, 3.3~V logic); LCMU is based on the Adafruit
Feather~RP2040 Adalogger (dual-core Cortex-M0+, 133~MHz). All four run
firmware written in C++17 using the Arduino / PlatformIO framework.

%═══════════════════════════════════════════════════════════════════════════════
\section{Propellant System}
\label{sec:propellant}
%═══════════════════════════════════════════════════════════════════════════════

\subsection{Flow Architecture}

Figure~\ref{fig:propellant} shows the propellant system schematic.

\begin{figure}[H]
  \centering
  \resizebox{0.82\textwidth}{!}{% propellant_flow.tex
% Propellant system schematic for inclusion in KJO_System_Overview.tex.
% \input this file inside a figure environment.
% Requires: KJO_tikz_styles.tex already loaded.

\begin{tikzpicture}[x=1cm, y=1cm]

%────────────────────────────────────────────────────────────────────
% Region boundaries
%────────────────────────────────────────────────────────────────────
% GSE boundary (dashed green)
\draw[gseboundary] (-1.0, 6.2) rectangle (5.0, 13.8);
\node[font=\footnotesize\bfseries, text=KJOgreen!60!black, anchor=north west]
  at (-1.0, 13.8) {Ground Support Equipment (GSE)};

% Rocket boundary (solid red/orange)
\draw[rocketboundary] (-4.5, -9.2) rectangle (11.5, 6.0);
\node[font=\footnotesize\bfseries, text=KJOred!60!black, anchor=south east]
  at (11.5, 6.0) {Rocket / Flight Hardware};

% Umbilical label
\node[font=\footnotesize\itshape, text=KJOgray] at (2.0, 6.5)
  {\textemdash\;\textit{Umbilical Interface}\;\textemdash};

%────────────────────────────────────────────────────────────────────
% GSE FILL PATH (left column, x=2)
%────────────────────────────────────────────────────────────────────
% N2O source tank
\node[srctank, minimum height=2.2cm] (src)
  at (2, 12.3) {N$_2$O\\Source\\Tank};

% Fill Valve (GSEMU-controlled servo)
\node[gseval, minimum width=2.4cm] (fillv)
  at (2, 9.8) {Fill Valve};
\node[font=\tiny, text=KJOgreen!70!black, above right=-1pt and 2pt of fillv]
  {GSEMU servo};

% QR Servo (GSEMU — umbilical release actuator)
\node[gseval, minimum width=2.4cm] (qrservo)
  at (2, 8.2) {QR Servo / Release};
\node[font=\tiny, text=KJOgreen!70!black, above right=-1pt and 2pt of qrservo]
  {GSEMU servo};

% Pipes: source → fill valve → QR servo
\draw[gsepipe] (src.south)    -- (fillv.north);
\draw[gsepipe] (fillv.south)  -- (qrservo.north);
% Arrow into rocket (downward through umbilical)
\draw[gsepipearrow] (qrservo.south) -- (2, 6.8);

%────────────────────────────────────────────────────────────────────
% ROCKET FILL PATH continues
%────────────────────────────────────────────────────────────────────
% QR Receiver (rocket side of umbilical)
\node[component] (qrrcvr) at (2, 5.5) {QR Receiver};

% Backflow prevention (check valve)
\node[component] (chkv)   at (2, 4.2) {Check Valve};

% T-junction dot
\coordinate (tjunc) at (2, 3.0);
\fill (tjunc) circle (3pt);

% Pipes: QR receiver → check valve → T-junction
\draw[pipe] (2, 6.8)     -- (qrrcvr.north);
\draw[pipe] (qrrcvr.south) -- (chkv.north);
\draw[pipe] (chkv.south)   -- (tjunc);

%────────────────────────────────────────────────────────────────────
% DUMP BRANCH (horizontal, left of T-junction)
%────────────────────────────────────────────────────────────────────
\node[emuval, minimum width=2.2cm] (dumpv) at (-2.0, 3.0) {Dump Valve};
\node[font=\tiny, text=KJOred, below=1pt of dumpv] {EMU servo};

% Vent arrow
\draw[-{Stealth[length=7pt,width=5pt]}, line width=2pt, draw=KJOgray!70]
  (dumpv.west) -- +(-1.5, 0)
  node[anchor=east, font=\footnotesize, text=KJOgray] {Vent};

% Pipe: T-junction → dump valve
\draw[pipe] (tjunc) -- (dumpv.east);

%────────────────────────────────────────────────────────────────────
% N2O OXIDIZER TANK (rocket)
%────────────────────────────────────────────────────────────────────
\node[n2otank, minimum height=2.8cm] (n2otank) at (2, 0.9) {N$_2$O\\Oxidizer\\Tank};

% Pipe: T-junction → N2O tank
\draw[pipe] (tjunc) -- (n2otank.north);

% Pressure transducer on N2O tank
\node[instrument] (pt1) at (4.8, 1.6) {PT};
\draw[sensorlead] (n2otank.east |- pt1) -- (pt1.west);
\node[font=\tiny, anchor=west, text=KJOgray] at (5.15, 1.6) {Tank Pressure};

% Temperature transducer on N2O tank
\node[instrument] (tt1) at (4.8, 0.3) {TT};
\draw[sensorlead] (n2otank.east |- tt1) -- (tt1.west);
\node[font=\tiny, anchor=west, text=KJOgray] at (5.15, 0.3) {Tank Temp};

%────────────────────────────────────────────────────────────────────
% PRESSURIZATION LINE (N2O vapor → Fuel Tank)
%────────────────────────────────────────────────────────────────────
% Horizontal dashed orange arrow from N2O tank to fuel tank at y=0.9
\draw[presspipearrow]
  (n2otank.east |- 0, 0.9) -- (6.2, 0.9)
  node[midway, above, font=\tiny, text=KJOorange] {Autogenous Press.};

%────────────────────────────────────────────────────────────────────
% OX FEED PATH
%────────────────────────────────────────────────────────────────────
\node[emuval, minimum width=2.4cm] (oxv) at (2, -1.8) {Ox Feed Valve};
\node[font=\tiny, text=KJOred, above right=-1pt and 2pt of oxv]
  {EMU servo, proportional};

\draw[oxpipe] (n2otank.south) -- (oxv.north);
% Arrow into chamber inlet
\draw[oxpipearrow] (oxv.south) -- (2, -4.8);

%────────────────────────────────────────────────────────────────────
% FUEL TANK AND FUEL PATH (right column, x=8)
%────────────────────────────────────────────────────────────────────
\node[fueltank, minimum height=2.2cm] (fueltank) at (8, 0.9) {Fuel\\Tank};

% Pressurization line arrives from N2O tank
\draw[presspipe] (6.2, 0.9) -- (fueltank.west);

% Fuel Feed Valve
\node[emuval, minimum width=2.4cm] (fuelv) at (8, -1.8) {Fuel Feed Valve};
\node[font=\tiny, text=KJOred, above right=-1pt and 2pt of fuelv]
  {EMU servo, proportional};

% Regenerative Cooling Jacket
\node[component, minimum width=2.6cm] (jacket) at (8, -3.4) {Regen.\ Cooling Jacket};

% Injector Plate
\node[component, minimum width=2.6cm] (injplate) at (8, -4.8) {Injector Plate\\(Fuel Injectors)};

% Fuel path pipes
\draw[fuelpipe] (fueltank.south) -- (fuelv.north);
\draw[fuelpipe] (fuelv.south)    -- (jacket.north);
\draw[fuelpipe] (jacket.south)   -- (injplate.north);
\draw[fuelpipearrow] (injplate.south) -- (8, -6.0);

%────────────────────────────────────────────────────────────────────
% COMBUSTION CHAMBER
%────────────────────────────────────────────────────────────────────
\node[rectangle, draw=black, line width=2pt, fill=yellow!10,
      minimum width=9.0cm, minimum height=1.3cm,
      font=\normalsize\bfseries, align=center]
  (chamber) at (5.0, -6.65) {Combustion Chamber};

% Ox enters from top-left
\draw[oxpipearrow] (2, -4.8) -- (2, -6.0);

% Fuel/injector enters from top-right
\draw[fuelpipearrow] (8, -6.0) -- (8, -6.0 |- chamber.north);

% Chamber pressure transducer
\node[instrument] (pt2) at (10.8, -6.65) {PT};
\draw[sensorlead] (chamber.east) -- (pt2.west);
\node[font=\tiny, anchor=west, text=KJOgray] at (11.15, -6.65) {Chamber\\Pressure};

% Thermocouple(s) on chamber wall
\node[instrument] (tc1) at (3.5, -5.85) {TC};
\draw[sensorlead] (tc1.south) -- (3.5, -6.0 |- chamber.north);
\node[font=\tiny, anchor=south, text=KJOgray] at (3.5, -5.55) {TC1};

\node[instrument] (tc2) at (6.5, -5.85) {TC};
\draw[sensorlead] (tc2.south) -- (6.5, -6.0 |- chamber.north);
\node[font=\tiny, anchor=south, text=KJOgray] at (6.5, -5.55) {TC2};

%────────────────────────────────────────────────────────────────────
% NOZZLE
%────────────────────────────────────────────────────────────────────
% Nozzle drawn as a trapezoid: wide at top (chamber exit), narrow at throat
\draw[line width=2pt]
  (3.0, -7.3) -- (2.5, -9.0)
  (7.0, -7.3) -- (7.5, -9.0)
  (2.5, -9.0) -- (7.5, -9.0);
\draw[line width=1.5pt] (3.0, -7.3) -- (7.0, -7.3);
\node[font=\small\itshape] at (5.0, -8.2) {Nozzle};

% Exhaust arrow
\draw[-{Stealth[length=9pt,width=6pt]}, line width=3pt, draw=KJOgray!60]
  (5.0, -9.0) -- (5.0, -9.8)
  node[anchor=north, font=\footnotesize, text=KJOgray] {Exhaust};

%────────────────────────────────────────────────────────────────────
% LEGEND
%────────────────────────────────────────────────────────────────────
\begin{scope}[xshift=-4.2cm, yshift=-7.5cm]
  \draw[gsepipe, line width=1.5pt] (0, 0) -- (0.9, 0);
  \node[anchor=west, font=\tiny] at (0.95, 0) {N$_2$O Fill (GSE)};
  \draw[oxpipe,  line width=1.5pt] (0, -0.4) -- (0.9, -0.4);
  \node[anchor=west, font=\tiny] at (0.95, -0.4) {Oxidizer (N$_2$O)};
  \draw[fuelpipe, line width=1.5pt] (0, -0.8) -- (0.9, -0.8);
  \node[anchor=west, font=\tiny] at (0.95, -0.8) {Fuel};
  \draw[presspipe, line width=1.2pt] (0, -1.2) -- (0.9, -1.2);
  \node[anchor=west, font=\tiny] at (0.95, -1.2) {Pressurization};
\end{scope}

\end{tikzpicture}
}
  \caption{Propellant system schematic. Green lines: N\textsubscript{2}O fill path (GSE).
           Blue lines: oxidizer feed path. Red lines: fuel feed path.
           Orange dashed: autogenous pressurization.}
  \label{fig:propellant}
\end{figure}

\subsection{Fill Sequence}

\begin{enumerate}[nosep]
  \item Operator commands \cmd{BEGIN\_FILL} from the RCU.
  \item EMU asserts the CMD bit (AUX\_IO\_1) high.
  \item GSEMU detects the CMD rising edge and opens the Fill valve servo.
  \item N\textsubscript{2}O flows: GSE source tank~$\to$ Fill valve~$\to$
        QR servo~$\to$ umbilical~$\to$ check valve~$\to$ on-board
        oxidizer tank.
  \item GSEMU asserts STATUS (AUX\_IO\_2) high when fill is complete.
  \item EMU de-asserts CMD; GSEMU closes the Fill valve.
\end{enumerate}

\subsection{Engine Firing Sequence}

\begin{enumerate}[nosep]
  \item Operator commands \cmd{ENABLE\_LCO\_CONTROL} (enables launch authority).
  \item Operator commands \cmd{START\_ENGINE}, which calls
        \sw{Execute\_Command(START\_ENGINE)}.
  \item EMU simultaneously opens Ox Feed valve and Fuel Feed valve to
        programmed positions and activates the igniter relay.
  \item Proportional control commands (\cmd{OXIDIZER\_FLOW\_PROP},
        \cmd{FUEL\_FLOW\_PROP}) may be sent during burn to adjust mixture ratio.
  \item \cmd{STOP\_ENGINE\_FLOWS} closes both feed valves to terminate thrust.
\end{enumerate}

\subsection{Dump / Safe Procedure}

The \cmd{DUMP} command opens the Oxidizer Feed valve (EMU-controlled), venting
the on-board N\textsubscript{2}O tank pressurant through the combustion chamber.
There is no dedicated Dump valve.  The RCU Dump button has latching behaviour
--- a second press sends \cmd{STOP\_DUMP} and closes the Oxidizer Feed valve.

\subsection{Pressurization}

The on-board N\textsubscript{2}O tank is self-pressurizing (autogenous).  The
vapour-space pressure is communicated to the fuel tank via a dedicated
pressurization line, eliminating the need for a separate pressurant.

%═══════════════════════════════════════════════════════════════════════════════
\section{Electronic Control Architecture}
\label{sec:architecture}
%═══════════════════════════════════════════════════════════════════════════════

Figure~\ref{fig:control} shows the four-unit electronic control architecture.

\begin{figure}[H]
  \centering
  \resizebox{\textwidth}{!}{% control_architecture.tex — Electronic control system block diagram.
% \input inside a figure environment. Requires KJO_tikz_styles.tex.

\begin{tikzpicture}[x=1cm, y=1cm, node distance=1.2cm]

%── RCU (left) ────────────────────────────────────────────────────────────────
\node[rectangle, draw=KJOblue, line width=1.5pt, fill=KJOblue!8,
      rounded corners=6pt, minimum width=4.0cm, minimum height=5.5cm,
      font=\small\bfseries, align=center, label={[font=\small\bfseries,
      text=KJOblue]above:Remote Control Unit (RCU)}]
  (rcu) at (0, 0) {};

\node[mcublock, minimum width=3.5cm] (rcu_mcu) at (0, 1.8) {Feather M0};
\node[hwblock, minimum width=3.5cm] (rcu_disp) at (0, 0.7) {3.5\textquotedbl\ TFT\\+ Touchscreen};
\node[hwblock, minimum width=3.5cm] (rcu_rad) at (0, -0.4) {RFM95 LoRa};
\node[hwblock, minimum width=3.5cm] (rcu_sd) at (0, -1.5) {SD + RTC};
\node[powerblock, minimum width=3.5cm] (rcu_bat) at (0, -2.4) {LiPo Battery};

%── EMU (center) ──────────────────────────────────────────────────────────────
\node[rectangle, draw=KJOred, line width=1.5pt, fill=KJOred!5,
      rounded corners=6pt, minimum width=4.2cm, minimum height=8.0cm,
      font=\small\bfseries, align=center, label={[font=\small\bfseries,
      text=KJOred]above:Engine Management Unit (EMU)}]
  (emu) at (6.5, -1.3) {};

\node[mcublock, minimum width=3.7cm] (emu_mcu) at (6.5, 2.0) {Feather M0};
\node[hwblock, minimum width=3.7cm] (emu_rad) at (6.5, 1.0) {RFM95 LoRa};
\node[hwblock, minimum width=3.7cm] (emu_oled) at (6.5, 0.0) {SH1107 OLED};
\node[hwblock, minimum width=3.7cm] (emu_gpio) at (6.5, -1.0) {MCP23X17 GPIO\\Expander};
\node[hwblock, minimum width=3.7cm] (emu_pwm) at (6.5, -2.0) {PWM Servo Wing\\(Ox, Fuel)};
\node[hwblock, minimum width=3.7cm] (emu_adc) at (6.5, -3.0) {ADS1015 ADC\\+ MAX31856 TC};
\node[hwblock, minimum width=3.7cm] (emu_sd) at (6.5, -4.0) {SD + RTC};
\node[powerblock, minimum width=3.7cm] (emu_bat) at (6.5, -4.8) {LiPo Battery};

%── GSEMU (right of EMU) ──────────────────────────────────────────────────────
\node[rectangle, draw=KJOgreen, line width=1.5pt, fill=KJOgreen!5,
      rounded corners=6pt, minimum width=4.0cm, minimum height=5.5cm,
      font=\small\bfseries, align=center, label={[font=\small\bfseries,
      text=KJOgreen]above:GSE Mgmt Unit (GSEMU)}]
  (gsemu) at (13.0, 0) {};

\node[mcublock, minimum width=3.5cm] (gse_mcu) at (13.0, 1.8) {Feather M0};
\node[hwblock, minimum width=3.5cm] (gse_oled) at (13.0, 0.7) {SH1107 OLED};
\node[hwblock, minimum width=3.5cm] (gse_pwm) at (13.0, -0.4) {PWM Servo Wing\\(Fill Valve)};
\node[hwblock, minimum width=3.5cm] (gse_sd) at (13.0, -1.5) {SD + RTC};
\node[powerblock, minimum width=3.5cm] (gse_bat) at (13.0, -2.4) {LiPo Battery};

%── LCMU (right of GSEMU) ─────────────────────────────────────────────────────
\node[rectangle, draw=KJOpurple, line width=1.5pt, fill=KJOpurple!5,
      rounded corners=6pt, minimum width=4.0cm, minimum height=5.5cm,
      font=\small\bfseries, align=center, label={[font=\small\bfseries,
      text=KJOpurple]above:Load Cell Mgmt Unit (LCMU)}]
  (lcmu) at (19.5, 0) {};

\node[mcublock, minimum width=3.5cm] (lcmu_mcu)  at (19.5, 1.8) {Feather RP2040\\Adalogger};
\node[hwblock, minimum width=3.5cm]  (lcmu_oled) at (19.5, 0.6) {SH1107 OLED};
\node[hwblock, minimum width=3.5cm]  (lcmu_cell) at (19.5, -0.5) {Mux + 3$\times$NAU7802\\(Load Cells)};
\node[hwblock, minimum width=3.5cm]  (lcmu_sd)   at (19.5, -1.6) {SD + RTC};
\node[powerblock, minimum width=3.5cm] (lcmu_bat) at (19.5, -2.4) {LiPo Battery};

%── LoRa Link (RCU ↔ EMU) ─────────────────────────────────────────────────────
\draw[radioline]
  (rcu.east |- rcu_rad) -- (emu.west |- emu_rad)
  node[midway, above, font=\footnotesize\itshape, text=KJOblue!70]
    {900 MHz LoRa};
\node[font=\tiny, text=KJOgray, align=center]
  at (3.25, 0.55) {Commands\\Responses};

%── CAN Bus (EMU -- GSEMU -- LCMU) ────────────────────────────────────────────
% Shared bus, not point-to-point: replaces the retired CMD/STATUS GPIO
% handshake for Fill-valve coordination, and is LCMU's only link to the
% other two units. Drawn as one bus line with tap stubs, matching CAN's
% actual multi-drop topology rather than the old wired handshake's
% point-to-point arrows.
% Routed below all three outer boxes (clear of every inner node) rather
% than through their interior -- GSEMU/LCMU tap in via a short vertical
% stub from their box's south edge; EMU's box already extends past this
% height, so its east-edge border point is used directly, no stub needed.
\coordinate (canbus_y) at (0, -3.3);
\draw[busline, -, line width=2pt]
  (emu.east |- canbus_y) -- (lcmu.east |- canbus_y)
  node[midway, below, font=\tiny, text=KJOgray] {CAN bus (1 Mbps)};
\draw[busline, -] (gsemu.south) -- (gsemu.south |- canbus_y);
\draw[busline, -] (lcmu.south)  -- (lcmu.south  |- canbus_y);

%── Umbilical disconnect-sense (EMU <-> GSEMU only, separate from CAN) ────────
\draw[gpioline, draw=KJOgray!70, -, dashed]
  ([yshift=-14pt] emu.east |- emu_gpio) --
  ([yshift=-14pt] gsemu.west |- emu_gpio)
  node[midway, below, font=\tiny, text=KJOgray!70!black] {disconnect sense (ground-ref.)};

%── Physical separation note ──────────────────────────────────────────────────
\node[font=\footnotesize\itshape, text=KJOgray, align=center]
  at (3.25, -5.2) {Operator\\(remote)};
\node[font=\footnotesize\itshape, text=KJOgray, align=center]
  at (6.5, -5.85) {In-rocket};
\node[font=\footnotesize\itshape, text=KJOgray, align=center]
  at (13.0, -4.3) {Launch pad\\(ground)};
\node[font=\footnotesize\itshape, text=KJOgray, align=center]
  at (19.5, -4.3) {Test stand};

%── Ground line ───────────────────────────────────────────────────────────────
\draw[KJOgray!40, line width=0.8pt, dashed]
  (0, -5.0) -- (6.5, -5.0)
  node[midway, below, font=\tiny, text=KJOgray] {radio link};

\end{tikzpicture}
}
  \caption{Electronic control architecture. The RCU communicates with the EMU
           over a 900~MHz LoRa radio link. EMU, GSEMU, and LCMU share a CAN
           bus (1~Mbps); a separate ground-referenced line gives EMU and
           GSEMU instant umbilical disconnect-sense, independent of CAN.}
  \label{fig:control}
\end{figure}

\subsection{RCU--EMU Radio Link}

The link uses the HopeRF RFM95W module (Semtech SX1276 chipset) operating at
915.0~MHz with 23~dBm transmit power.  Packets are simple binary structures
(see Section~\ref{sec:packets}).  The RCU sends a \cmd{HEARTBEAT} packet at
5-second intervals when no other command is pending; the EMU echoes a
\cmd{HEARTBEAT} response containing valve positions and battery voltage.  Link
loss is declared after 30~seconds without a response.

\subsection{EMU--GSEMU--LCMU CAN Bus}

EMU, GSEMU, and LCMU share one CAN bus (MCP2515 controllers, 1~Mbps,
node IDs EMU~=~1, GSEMU~=~2, LCMU~=~3) --- addressed command/response only,
no unsolicited broadcasts. This replaced an earlier wired GPIO CMD/STATUS
handshake that carried Fill-valve coordination directly between EMU and
GSEMU; that handshake's two lines (\sig{AUX\_IO\_1}/\sig{AUX\_IO\_2} on
EMU's MCP23X17 GPIO expander) are retired. Full command tables and protocol
details for each unit are in their own reference documents
(Section~\ref{sec:references}).

Separately, a ground-referenced disconnect-sense line on each side of the
EMU--GSEMU umbilical (not CAN) gives both units instant physical-disconnect
detection independent of bus traffic or timeout logic.

%═══════════════════════════════════════════════════════════════════════════════
\section{Shared Libraries}
\label{sec:libs}
%═══════════════════════════════════════════════════════════════════════════════

Common definitions reside in \texttt{KJO\_Shared\_Libraries} and are shared
across all four firmware projects via PlatformIO library dependencies:

\begin{center}
\begin{tabularx}{\textwidth}{@{}lYl@{}}
\toprule
Library & Contents & Used by \\
\midrule
\sw{KJO\_System\_Config}  & Platform IDs, version strings      & RCU, EMU, GSEMU, LCMU \\
\sw{KJO\_Command\_Defs}   & LoRa command codes, packet structs  & RCU, EMU \\
\sw{KJO\_CAN\_Command\_Defs} & CAN command codes, packet structs, node IDs & EMU, GSEMU, LCMU \\
\sw{KJO\_Radio}           & Radio pin assignments, frequency    & RCU, EMU \\
\sw{KJO\_RCU\_Display}    & TFT GUI functions and data types    & RCU \\
\bottomrule
\end{tabularx}
\end{center}

%═══════════════════════════════════════════════════════════════════════════════
\section{Command Set}
\label{sec:commands}
%═══════════════════════════════════════════════════════════════════════════════

\subsection{Packet Structures}

Three binary packet types are used on the radio link:

\begin{itemize}[nosep]
  \item \textbf{Command\_Payload} (4 bytes): \texttt{command} (2B),
        \texttt{param} (short, 2B).  Used for most commands.
  \item \textbf{Heartbeat\_Command\_Packet} (6 bytes): \texttt{command} (2B),
        \texttt{RCU\_voltage} (float, 4B).  Sent by RCU as heartbeat.
  \item \textbf{Heartbeat\_Response\_Packet} (10 bytes): \texttt{command} (2B),
        \texttt{EMU\_voltage} (float, 4B),
        \texttt{fill\_pct}, \texttt{ox\_pct}, \texttt{fuel\_pct},
        \texttt{dump\_pct} (uint8, 1B each).  Sent by EMU.
  \item \textbf{Command\_Return\_Packet}: \texttt{command} (2B),
        \texttt{status} (bool), \texttt{r\_val} (short).
        Returned for non-heartbeat commands.
\end{itemize}

\subsection{Command Table}

\begingroup
\footnotesize
\begin{longtable}{@{}rlp{2.1cm}p{1.3cm}p{5.0cm}@{}}
\toprule
Code & Name & Direction & \texttt{param} & Description \\
\midrule
\endhead
\bottomrule
\endfoot
 1 & \cmd{BEGIN\_FILL}           & RCU$\to$EMU & ---   & Command GSEMU to open Fill valve \\
 2 & \cmd{STOP\_FILL}            & RCU$\to$EMU & ---   & Command GSEMU to close Fill valve \\
 3 & \cmd{OPEN\_OXIDIZER\_FEED}  & RCU$\to$EMU & ---   & Open Ox Feed valve fully \\
 4 & \cmd{CLOSE\_OXIDIZER\_FEED} & RCU$\to$EMU & ---   & Close Ox Feed valve \\
 5 & \cmd{OPEN\_FUEL\_FEED}      & RCU$\to$EMU & ---   & Open Fuel Feed valve fully \\
 6 & \cmd{CLOSE\_FUEL\_FEED}     & RCU$\to$EMU & ---   & Close Fuel Feed valve \\
 7 & \cmd{OXIDIZER\_FLOW\_PROP}  & RCU$\to$EMU & \%    & Set Ox valve position (0--100\%) \\
 8 & \cmd{FUEL\_FLOW\_PROP}      & RCU$\to$EMU & \%    & Set Fuel valve position (0--100\%) \\
 9 & \cmd{DUMP}                  & RCU$\to$EMU & ---   & Open Ox valve; vent pressurant through chamber \\
10 & \cmd{STOP\_DUMP}            & RCU$\to$EMU & ---   & Close Ox valve \\
11 & \cmd{START\_ENGINE\_FLOWS}  & RCU$\to$EMU & ---   & Open both feed valves simultaneously \\
12 & \cmd{STOP\_ENGINE\_FLOWS}   & RCU$\to$EMU & ---   & Close both feed valves simultaneously \\
13 & \cmd{ENABLE\_LCO\_CONTROL}  & RCU$\to$EMU & ---   & Enable launch authority \\
14 & \cmd{START\_ENGINE}         & RCU$\to$EMU & ---   & Fire igniter and open valves \\
15 & \cmd{RESET}                 & RCU$\to$EMU & ---   & Reset EMU to safe state \\
16 & \cmd{READ\_TANK\_PRESSURE}  & RCU$\to$EMU & ---   & Return oxidizer tank pressure \\
17 & \cmd{READ\_CHAMBER\_PRESSURE} & RCU$\to$EMU & ---  & Return chamber pressure \\
18 & \cmd{READ\_TANK\_TEMP}      & RCU$\to$EMU & ---   & Return oxidizer tank temperature \\
19 & \cmd{READ\_CHAMBER\_TEMP}   & RCU$\to$EMU & ---   & Return combustion chamber temperature \\
20 & \cmd{START\_PRESSURE\_RECORDING} & RCU$\to$EMU & --- & Begin high-rate pressure logging \\
21 & \cmd{STOP\_PRESSURE\_RECORDING}  & RCU$\to$EMU & --- & End high-rate pressure logging \\
22 & \cmd{READ\_VALVE\_POSITION} & RCU$\to$EMU & valve\# & Return servo position (0--100\%) \\
23 & \cmd{HEARTBEAT}             & Bidirectional & ---  & Keepalive / link maintenance \\
24 & \cmd{RCU\_STATUS}           & RCU$\to$EMU & ---   & Request full status packet \\
\end{longtable}
\endgroup

%═══════════════════════════════════════════════════════════════════════════════
\section{Document References}
\label{sec:references}
%═══════════════════════════════════════════════════════════════════════════════

\begin{center}
\begin{tabularx}{\textwidth}{@{}llY@{}}
\toprule
Document & File & Contents \\
\midrule
RCU Reference   & \texttt{RCU/KJO\_RCU\_Reference.pdf}   & Hardware, software, GUI \\
EMU Reference   & \texttt{EMU/KJO\_EMU\_Reference.pdf}   & Hardware, software, valve control \\
GSEMU Reference & \texttt{GSEMU/KJO\_GSEMU\_Reference.pdf} & Hardware, software, Fill valve \\
LCMU Reference  & \texttt{LCMU/KJO\_LCMU\_Reference.pdf} & Hardware, software, CAN bus, load cells \\
\bottomrule
\end{tabularx}
\end{center}

%═══════════════════════════════════════════════════════════════════════════════
\section{Version History}
%═══════════════════════════════════════════════════════════════════════════════

\begin{center}
\begin{tabular}{@{}lll@{}}
\toprule
Version & Date & Notes \\
\midrule
0.1 & March 2026 & Initial release. \\
\bottomrule
\end{tabular}
\end{center}

\end{document}
